\documentclass[11pt,a4paper]{article}
\usepackage[T1]{fontenc}
\usepackage{lmodern}
\usepackage[margin=1in]{geometry}
\usepackage{amsmath,amssymb,amsthm,mathtools}
\usepackage{booktabs,microtype,enumitem,needspace,xcolor}
\usepackage{hyperref}
\hypersetup{colorlinks=true,linkcolor=blue!45!black,citecolor=blue!45!black,
 urlcolor=blue!45!black,
 pdftitle={A positive constant-recursive sequence with nonquasipolynomial power ranks},
 pdfauthor={Dev D. Goyal},
 pdfsubject={Exact minimal recurrences for termwise powers}}
\setlength{\parindent}{0pt}
\setlength{\parskip}{0.35em}
\newtheorem{theorem}{Theorem}
\newtheorem{lemma}[theorem]{Lemma}
\newtheorem{corollary}[theorem]{Corollary}
\newcommand{\rank}{\operatorname{rank}}
\newcommand{\NN}{\mathbb N}
\newcommand{\CC}{\mathbb C}
\newcommand{\QQ}{\mathbb Q}
\newcommand{\ZZ}{\mathbb Z}
\newcommand{\coeff}[2]{[#1]#2}
\title{A positive constant-recursive sequence\\with nonquasipolynomial power ranks}
\author{Dev D. Goyal}
\date{28 October 2026; priority note added 30 October 2026}
\begin{document}
\maketitle
\begin{quote}\small
\textbf{Autonomous AI research disclosure.}
This work is fully autonomous agentic research. GPT 6.1 Sol in Codex
produced formal proofs and verification software for the accompanying
package. Dev D. Goyal is responsible for publication.
\end{quote}
\begin{abstract}
Termwise powers of a constant-recursive sequence again satisfy constant-coefficient
recurrences, but the minimal orders can depend on cancellations in the power
expansion. An explicit positive integer sequence of rank thirteen has power ranks
equal to a cubic polynomial minus twice a parity-restricted square count.
The correction is $(\sqrt3-1)\sqrt M+O(1)$ at power $M$, so these ranks are
not eventually quasipolynomial. The proof determines every characteristic-root
multiplicity for every positive power. A companion construction gives infinitely
many distinct power-rank sequences at each fixed initial rank $r\ge3$.
The recurrence conventions and the standard algebraic tools are stated explicitly.
\end{abstract}

\section*{Priority and subsequent correspondence}
\emph{Added 30 October 2026.} On 29 October 2026, following my correspondence with Professor Eric Rowland, he informed me that he had recently been experimenting with an LLM and had independently obtained an infinite family of rank sequences for rank-three recurrences. That family also challenges the fixed-rank finiteness assertion of~\cite{RowlandSistos} discussed at the end of this paper. I had not known of this work, and the novelty searches carried out by my research agents had not found it. I acknowledge this independently discovered, overlapping result. I do not suggest that his construction duplicates the rank-thirteen example or the other companion theorems; the constructions have not been compared. The proofs and the formalization are unchanged.

\section{The question and the main result}
Taking a termwise power of a sequence can increase the order of its shortest
constant-coefficient recurrence. An exponential expansion explains this increase:
products of the original exponential bases become new bases. Equal products
must be grouped, and their polynomial coefficients can cancel. The actual
recurrence rank depends on the resulting nonzero coefficients and their degrees.
For related work on ranks of powers, see Rowland and Sistos Barr\'on~\cite{RowlandSistos}.

The main example below has a different cancellation pattern at square indices.
These losses are small compared with the total rank, but they prevent every
eventual quasipolynomial description. The characteristic polynomial is determined
explicitly, so a loss in one coefficient is not mistaken for a larger hidden loss.

\paragraph{Recurrence conventions.}
All sequences are indexed by $n\in\NN=\{0,1,2,\ldots\}$.
Power exponents $M$ are nonnegative integers.
Recurrence coefficients are initially allowed in $\CC$.
Let $E$ be the forward shift, $(Eu)(n)=u(n+1)$.
A sequence is \emph{constant-recursive} if some monic polynomial $q\in\CC[T]$
with $q(0)\ne0$ satisfies $q(E)u(n)=0$ for all sufficiently large $n$.
Its \emph{rank} is the least degree of such a polynomial, and its monic
minimal polynomial is its \emph{minimal characteristic polynomial}.
An eventually zero sequence has rank zero and polynomial $1$ by convention.
For positive order $R$, the characteristic polynomial
\[
 T^R+a_{R-1}T^{R-1}+\cdots+a_0,\qquad a_0\ne0,
\]
gives the eventual recurrence
\[
 u(n+R)=-\sum_{h=0}^{R-1}a_hu(n+h).
\]
All sequences constructed here are everywhere positive, and their displayed
recurrences hold for every $n\ge0$. Their minimal polynomials lie in $\ZZ[T]$.
Consequently their ranks are the same over $\QQ$, $\mathbb R$, and $\CC$.

For $M\ge0$, the power $u^M$ means the termwise power $u^M(n)=u(n)^M$.
The sequences here never vanish, so $u^0$ is the constant sequence one and has rank one.
A function of $M$ is \emph{eventually quasipolynomial} if there are an integer
$\ell\ge1$, a threshold $M_0$, and real polynomials $f_0,\ldots,f_{\ell-1}$
such that its value is $f_j(M)$ whenever $M\ge M_0$ and $M\equiv j\pmod\ell$.

Define
\begin{equation}\label{eq:s}
 s(n)=1+2^n+3^n(1-2^n)^2+n9^n(1+2^n)^3.
\end{equation}
For $M\ge1$, put
\begin{equation}\label{eq:Q}
 Q(M)=\#\{q\in\NN:M<q^2<3M,\ q^2-M\text{ is odd}\}.
\end{equation}

\begin{theorem}\label{thm:main}
The sequence $s$ in~\eqref{eq:s} is positive integer-valued and has rank $13$.
For every $M\ge1$, its $M$th termwise power has minimal characteristic polynomial
\begin{equation}\label{eq:chi}
 \chi_M(T)=\prod_{j=0}^{2M}\prod_{k=0}^{M+j}
 (T-2^k3^j)^{1+\lfloor j/2\rfloor-\epsilon_{M,j,k}},
\end{equation}
where
\[
 \epsilon_{M,j,k}=
 \begin{cases}
 1,&j\text{ odd and }(M+j-2k)^2=M+j,\\
 0,&\text{otherwise}.
 \end{cases}
\]
Factors with exponent zero are omitted. In particular,
\begin{equation}\label{eq:rank}
 R(M):=\rank(s^M)=P(M)-2Q(M),\qquad
 P(M)=\frac{(M+1)(2M+1)(7M+6)}6.
\end{equation}
For every $M\ge1$,
\begin{equation}\label{eq:asymp}
 \left|P(M)-R(M)-(\sqrt3-1)\sqrt M\right|\le2.
\end{equation}
The rank sequence $R$ is not eventually quasipolynomial, even with real
coefficients allowed in its constituent polynomials.
\end{theorem}

The first eight ranks, starting at $M=1$, are
\[
 13,\ 50,\ 124,\ 253,\ 451,\ 726,\ 1098,\ 1579.
\]
The minimal polynomial of $s$ itself is
\begin{equation}\label{eq:initialchi}
 \begin{split}
 &(T-1)(T-2)(T-3)(T-6)(T-12)\\[-2pt]
 &\hspace{1em}\cdot(T-9)^2(T-18)^2(T-36)^2(T-72)^2.
 \end{split}
\end{equation}
Thus the example has repeated characteristic roots. No minimality claim is made
for the number $13$ among examples with nonquasipolynomial power ranks.

\section{The rank principle}
The following standard exponential-polynomial fact supplies minimality, rather
than only an upper bound on recurrence order. A proof is included to cover the
eventual-recurrence convention; compare Theorem~2.1 of~\cite{RowlandSistos}.

\begin{lemma}\label{lem:rank}
Let $h\ge1$. Let $\lambda_1,\ldots,\lambda_h$ be distinct nonzero complex
numbers, and let $P_i\in\CC[n]$ be nonzero polynomials of degrees $e_i$.
For
\[
 u(n)=\sum_{i=1}^hP_i(n)\lambda_i^n,
\]
the minimal characteristic polynomial and rank are
\[
 \chi_u(T)=\prod_{i=1}^h(T-\lambda_i)^{e_i+1},\qquad
 \rank(u)=\sum_{i=1}^h(e_i+1).
\]
\end{lemma}
\begin{proof}
On a single component,
\[
 (E-\lambda)\bigl(P(n)\lambda^n\bigr)
 =\lambda^{n+1}\bigl(P(n+1)-P(n)\bigr).
\]
This lowers a positive polynomial degree by one and kills a constant.
If $\mu\ne\lambda$, then
\[
 (E-\mu)\bigl(P(n)\lambda^n\bigr)
 =\lambda^n\bigl(\lambda P(n+1)-\mu P(n)\bigr),
\]
which preserves the degree and multiplies its leading coefficient by
$\lambda-\mu$.

The distinct components are independent even on a tail. Indeed, if their sum
were eventually zero, apply $\prod_{j\ne i}(E-\lambda_j)^{e_j+1}$ for a fixed $i$.
All other components vanish. The remaining component is a nonzero polynomial
times $\lambda_i^n$, since the applied factors preserve its degree.
It cannot vanish at all sufficiently large integers.

Now let $q$ be a nonzero polynomial such that $q(E)$ annihilates the sum eventually. Write
$q(T)=(T-\lambda_i)^{a_i}g_i(T)$ with $g_i(\lambda_i)\ne0$.
The leading polynomial coefficient after $g_i(E)$ is multiplied by
$g_i(\lambda_i)$, so its degree remains $e_i$.
Applying $(E-\lambda_i)^{a_i}$ kills it exactly when $a_i\ge e_i+1$.
The independence just proved forces this inequality for every i.
Thus the displayed product divides every annihilator. Conversely, that product
annihilates all components for every $n\ge0$ and has nonzero constant term.
The polynomial and rank assertions follow.
\end{proof}

If the polynomial in Lemma~\ref{lem:rank} has integer coefficients, it gives
a recurrence over $\QQ$ and $\mathbb R$. No smaller recurrence over either
field can exist, since it would also be a smaller complex recurrence.
This proves the field-independence assertion for the constructions below.

\section{Every characteristic multiplicity of the main example}
First expand~\eqref{eq:s} as
\begin{equation}\label{eq:expand-s}
 \begin{split}
 s(n)={}&1+2^n+3^n-2\,6^n+12^n\\
 &+n\bigl(9^n+3\,18^n+3\,36^n+72^n\bigr).
 \end{split}
\end{equation}
The bases are distinct. Lemma~\ref{lem:rank} gives rank $5+2\cdot4=13$
and~\eqref{eq:initialchi}. The expression~\eqref{eq:s} proves positivity and
integrality for every $n\ge0$.

\subsection{A layer expansion}
Set
\[
 F(X,Y,Z)=(1+X)+Y(1-X)^2+ZY^2(1+X)^3.
\]
Then $s(n)=F(2^n,3^n,n)$. Define the polynomial
\[
 C_{M,j,k}(Z)=[X^kY^j]F(X,Y,Z)^M.
\]
After substitution, it is the coefficient polynomial of the exponential
base $2^k3^j$. Unique factorization makes these bases distinct for different
pairs $(j,k)$. Each nonzero $C_{M,j,k}$ contributes
$1+\deg C_{M,j,k}$ to the rank.

Choose $c$ factors from the Z-term, $j-2c$ from the middle term, and $M-j+c$
from the first term. This gives the exact expansion
\begin{equation}\label{eq:layers}
 [Y^j]F^M=
 \sum_{c=\max(0,j-M)}^{\lfloor j/2\rfloor}
 \frac{M!\,Z^c(1+X)^{M-j+4c}(1-X)^{2j-4c}}
 {(M-j+c)!(j-2c)!c!}.
\end{equation}
The range is exactly the condition that all three chosen counts be
nonnegative. Each multinomial scalar is positive. The $X$-degree of each
summand is $M+j$. Hence only $0\le j\le2M$ and $0\le k\le M+j$ occur.
For fixed $j$, different $c$ give different Z-degrees. There is only one
summand at each possible Z-degree, so its X-coefficient determines whether
that degree survives.

If $j=2a$, the top Z-degree is a. Its coefficient is a positive scalar
times $(1+X)^{M+2a}$. All coefficients for $0\le k\le M+2a$ are positive.
Thus every position in an even layer has degree exactly a and contributes
multiplicity $a+1$.

If $j=2a+1$, put $N=M+j$. The top Z-degree is a, with coefficient a
positive scalar times
\begin{equation}\label{eq:top}
 (1+X)^{N-2}(1-X)^2.
\end{equation}
Here $N\ge2$. The next lemma identifies all zeros of this coefficient and
shows what happens at the next Z-degree.

\subsection{Two coefficient identities}
These are standard binomial coefficient, or Krawtchouk, identities.
The derivation fixes the normalization needed for the degree calculation.

\begin{lemma}\label{lem:coeff}
For integers $N\ge2$ and $0\le k\le N$,
\begin{equation}\label{eq:quadratic}
 [X^k](1+X)^{N-2}(1-X)^2
 =\binom Nk\frac{(N-2k)^2-N}{N(N-1)}.
\end{equation}
If $N\ge6$ and $(N-2k)^2=N$, then
\begin{equation}\label{eq:sixth}
 [X^k](1+X)^{N-6}(1-X)^6
 =\binom Nk\frac{16}{(N-2)(N-3)(N-5)}>0.
\end{equation}
\end{lemma}
\begin{proof}
Let $t=N-2k$ and
\[
 K_r=[z^r](1-z)^k(1+z)^{N-k}.
\]
For $0\le r\le N$, expansion and cancellation of factorials give
\begin{equation}\label{eq:duality}
 [X^k](1-X)^r(1+X)^{N-r}
 =\frac{\binom Nk}{\binom Nr}K_r.
\end{equation}
Term by term, the needed identity is
\[
 \frac{\binom ri\binom{N-r}{k-i}}{\binom Nk}
 =\frac{\binom ki\binom{N-k}{r-i}}{\binom Nr},
\]
with out-of-range binomial coefficients zero.
The polynomial $G(z)=(1-z)^k(1+z)^{N-k}$ satisfies
$(1-z^2)G'(z)=(t-Nz)G(z)$. Comparing coefficients gives
\begin{equation}\label{eq:Krec}
 (r+1)K_{r+1}=tK_r-(N-r+1)K_{r-1},\qquad K_0=1,\quad K_1=t.
\end{equation}
In particular, $K_2=(t^2-N)/2$, proving~\eqref{eq:quadratic}.

Under $t^2=N$, recurrence~\eqref{eq:Krec} gives
\[
 \begin{aligned}
 K_2&=0,& K_3&=-\frac{t(N-1)}3,& K_4&=-\frac{N(N-1)}{12},\\
 K_5&=\frac{t(N-1)(N-4)}{20},&
 K_6&=\frac{N(N-1)(N-4)}{45}.
 \end{aligned}
\]
Since $\binom N6=N(N-1)(N-2)(N-3)(N-4)(N-5)/720$, equation~\eqref{eq:duality}
at $r=6$ proves~\eqref{eq:sixth}. Its denominator is positive for $N\ge6$.
\end{proof}

\subsection{Exactly one loss at every exceptional position}
Equation~\eqref{eq:quadratic} shows that the top coefficient in an odd
layer vanishes exactly when $N$ is a square, say $N=q^2$, at the two positions
\begin{equation}\label{eq:positions}
 k=\frac{N-q}{2},\qquad k=\frac{N+q}{2}.
\end{equation}
They are distinct because $N\ge2$ forces $q\ge2$. They are integers because
$q$ and $q^2$ have the same parity, and they lie in $[0,N]$.

If $a=0$, the coefficient polynomial has degree at most zero. A cancellation
therefore removes that base completely. If $a\ge1$, a lower degree must be
checked. The endpoint $a=M-1$ would give $N=3M-1\equiv2\pmod3$, impossible
for a square. Thus every exceptional layer with $a\ge1$ has $a\le M-2$.
It follows that $c=a-1$ is valid in~\eqref{eq:layers}: it is nonnegative and
$a-1\ge j-M=2a+1-M$.
The coefficient at that Z-degree is a positive scalar times
\[
 [X^k](1+X)^{N-6}(1-X)^6.
\]
Here $M\ge3$ and $N=M+2a+1\ge6$. Since $N$ is a square, $N\ge9$.
Equation~\eqref{eq:sixth} makes this coefficient strictly positive.
No other summand contributes to the same Z-degree. Thus the degree is
exactly $a-1$. In every exceptional position the characteristic multiplicity
drops by exactly one, and never by more.

For example, at $M=3,j=1$ the top factor is
$(1+X)^2(1-X)^2=1-2X^2+X^4$. There is no lower Z-degree, so the two bases
at $k=1,3$ disappear. At $M=6,j=3$, the exceptional $N$ is $9$.
The leading degree vanishes at $k=3,6$, while
$[X^3](1+X)^3(1-X)^6=8$. The next multinomial scalar is 20, so its
coefficient is 160, not zero.

Combining the even and odd layers with Lemma~\ref{lem:rank} proves the
minimal polynomial~\eqref{eq:chi}. Every factor has a nonnegative integer
exponent, and its base is a positive integer. Thus $\chi_M$ is monic,
integral, and has nonzero constant term. Its recurrence holds for every $n\ge0$.

\section{The square correction and its consequence}
The exceptional odd layers are $j=1,3,\ldots,2M-1$ with $M+j=q^2$.
They correspond exactly to~\eqref{eq:Q}. Each such layer loses two units
of characteristic multiplicity.
Before these losses, the rank sum is
\begin{align*}
 \sum_{j=0}^{2M}(M+j+1)(\lfloor j/2\rfloor+1)
 &=\sum_{a=0}^{M}(M+2a+1)(a+1)\\
 &\quad+\sum_{a=0}^{M-1}(M+2a+2)(a+1)\\
 &=\frac{(M+1)(M+2)(7M+3)}6
   +\frac{M(M+1)(7M+2)}6\\
 &=P(M).
\end{align*}
This proves~\eqref{eq:rank} for every $M\ge1$.

For an exact floor formula, let
\[
 L=\lfloor\sqrt M\rfloor,\qquad U=\lfloor\sqrt{3M-1}\rfloor,
 \qquad e=M\bmod2.
\]
The eligible $q$ satisfy $L<q\le U$ and have parity $1-e$. Therefore
\begin{equation}\label{eq:floor}
 Q(M)=\left\lfloor\frac{U+1-e}{2}\right\rfloor
      -\left\lfloor\frac{L+1-e}{2}\right\rfloor.
\end{equation}
This handles the strict square endpoints exactly.

The same $q$ lie in the open interval $(\sqrt M,\sqrt{3M})$ on one lattice
of spacing two. Its point count differs from half the interval length
by at most one. Hence
\[
 \left|Q(M)-\frac{\sqrt3-1}{2}\sqrt M\right|\le1,
\]
which proves~\eqref{eq:asymp}.

Suppose $R$ were eventually quasipolynomial with period $\ell\ge1$ and
real constituent polynomials. On each residue class modulo $\ell$, the
difference $P-R$ would eventually agree with a polynomial in $M$.
By~\eqref{eq:asymp}, that polynomial is unbounded and is $o(M)$ along the
class. A constant polynomial is bounded. A nonconstant polynomial cannot
be $o(M)$ along an unbounded positive arithmetic progression.
Both alternatives are impossible. This completes the proof of
Theorem~\ref{thm:main}.

\section{A companion family at every fixed rank}
This construction isolates a simpler mechanism: equal exponential products
without coefficient cancellation.

\begin{theorem}\label{thm:family}
Fix integers $r\ge3$, $d\ge2$, and distinct odd primes
$p_1,\ldots,p_{r-3}$, with an empty list when $r=3$. The positive integer sequence
\[
 u_{r,d}(n)=1+2^n+2^{dn}+\sum_{i=1}^{r-3}p_i^n
\]
has rank $r$. For every $M\ge0$,
\begin{equation}\label{eq:family-rank}
 \rank(u_{r,d}^M)=\binom{M+r-1}{r-1}
 -\begin{cases}
 0,&M<d,\\
 \binom{M-d+r-1}{r-1},&M\ge d.
 \end{cases}
\end{equation}
For each fixed $r\ge3$, varying d gives pairwise distinct power-rank sequences.
\end{theorem}
\begin{proof}
The $r$ original exponential bases are distinct, so Lemma~\ref{lem:rank}
gives rank $r$. Expand the Mth power. A term is indexed by a tuple of nonnegative integers
\[
 b=(b_0,b_1,b_2,c_1,\ldots,c_{r-3}),\qquad
 b_0+b_1+b_2+\sum_i c_i=M,
\]
and has base
\begin{equation}\label{eq:base}
 \lambda_b=2^{b_1+db_2}\prod_i p_i^{c_i}.
\end{equation}
All multinomial coefficients are positive, so equal bases do not cancel.

Write $b_1=qd+a$ with $0\le a<d$. Replace $b$ by
\[
 b^*=(b_0+(d-1)q,\ a,\ b_2+q,\ c_1,\ldots,c_{r-3}).
\]
All entries remain nonnegative. The coordinate sum and base are unchanged.
Thus every base has a representative with $b_1<d$.
Such a representative is unique: unique factorization determines the
odd-prime exponents and the exponent of two, Euclidean division by d
determines $b_1,b_2$, and the fixed sum determines $b_0$.

There are $\binom{M+r-1}{r-1}$ weak compositions of $M$ into $r$ parts.
When $M\ge d$, those with $b_1\ge d$ are counted by
$\binom{M-d+r-1}{r-1}$ after subtracting $d$ from $b_1$.
For $M<d$, there are none. Lemma~\ref{lem:rank} proves~\eqref{eq:family-rank}.

If $d<d^\prime$, compare the sequences at $M=d$. The first rank is
$\binom{d+r-1}{r-1}-1$, and the second is $\binom{d+r-1}{r-1}$.
Thus they differ.
\end{proof}

Let $\mathcal C_{r,d,M}$ be the weak compositions of M with $b_1<d$.
The proof also gives the explicit minimal characteristic polynomial
\begin{equation}\label{eq:family-chi}
 \chi_{r,d,M}(T)=\prod_{b\in\mathcal C_{r,d,M}}(T-\lambda_b).
\end{equation}
Every factor occurs once. This polynomial is monic, integral, and has
nonzero constant term. Together with~\eqref{eq:chi}, it specifies the
minimal recurrence of every power: if $\chi(T)=T^R+\sum_{h<R}a_hT^h$, then
\[
 u(n+R)^M=-\sum_{h=0}^{R-1}a_hu(n+h)^M.
\]
For $M=0$, both constructions give $\chi(T)=T-1$ and rank one.

The generating function for~\eqref{eq:family-rank} is
\[
 \sum_{M\ge0}\rank(u_{r,d}^M)z^M=\frac{1-z^d}{(1-z)^r}.
\]
For $M\ge d$, the rank is a polynomial of degree $r-2$ with leading coefficient
$d/(r-2)!$. When $r=3$, the formula becomes
\[
 \rank(u_{3,d}^M)=
 \begin{cases}
 \binom{M+2}{2},&0\le M\le d-1,\\[2pt]
 dM+1-\dfrac{(d-1)(d-2)}2,&M\ge d-1.
 \end{cases}
\]
Both branches agree at $M=d-1$. For example, $1+2^n+32^n$ has power ranks
\[
 3,6,10,15,20,25,30,\ldots
\]
and minimal recurrence
\[
 u(n+3)=35u(n+2)-98u(n+1)+64u(n).
\]

The generating function is a standard Hilbert-function calculation.
The graded monomial map
\[
 X_0\mapsto U,\quad X_1\mapsto UV,\quad X_2\mapsto UV^d,
 \quad X_{i+2}\mapsto UW_i
\]
has the relation $X_1^d=X_0^{d-1}X_2$.
The same normal form gives the monomial basis with exponent of $X_1$ less
than $d$, and hence the series $(1-z^d)/(1-z)^r$.
This is the familiar principal-relation mechanism, not a new Hilbert-series principle.

\section{Verification, attribution, and limits}
The arguments above prove the infinite statements. Exact finite checks
provide additional checks of the expansions and indexing.
The accompanying integer-arithmetic programs check canonical supports,
individual coefficient degrees, the floor formula, and characteristic
polynomials recovered from sequence values modulo primes.
One recorded run covered 715 support cases and 851,565 normal-form reductions,
20,000 floor-count identities, every coefficient multiplicity for $1\le M\le14$,
and direct recurrence comparisons for $M=1,2,3$ modulo both
$1{,}000{,}000{,}007$ and $1{,}000{,}000{,}009$.
These finite checks are corroboration, not a proof for all powers.

The accompanying Lean package formalizes Lemma~\ref{lem:rank} and
Theorem~\ref{thm:main}. It proves eventual recurrence minimality over
$\CC$, the full minimal polynomial for every power, the exact rank and
floor formulas, the error bound in~\eqref{eq:asymp}, and failure of eventual
quasipolynomiality with real constituent polynomials. It also proves
integrality of the minimal polynomials, equality of the ranks over
$\QQ$, $\mathbb R$, and $\CC$, and the recurrences from index zero.
The zeroth-power boundary is included.
The package also formalizes Theorem~\ref{thm:family} for every $r\ge3$,
$d\ge2$, and $M\ge0$. It proves the full companion minimal polynomial,
its integral recurrence from index zero, and equality of rational, real,
and complex ranks. It includes the existence of the required distinct
odd primes and infinitely many distinct power-rank sequences at each $r$.
The companion generating function, eventual polynomial degree and leading
coefficient, and simplified rank-three examples retain their conventional
proofs above and are outside the formal package.

The exponential-polynomial rank principle, binomial coefficient identities,
and Hilbert-function mechanism are standard supporting mathematics.
Rowland and Sistos Barr\'on~\cite[p.~2]{RowlandSistos} conjecture that, for each
$r\ge1$, only finitely many power-rank sequences arise from constant-recursive
integer sequences of rank $r$. Theorem~\ref{thm:family} contradicts this
unnumbered finiteness conjecture as written for every $r\ge3$; at rank three, Rowland has independently obtained an infinite family (see the priority note at the beginning of this paper). Its examples
are positive integer sequences with integral recurrences and positive simple
integer characteristic roots. The source's eventual-recurrence convention
agrees with the convention here. Its Theorem~5.4 concerns rank two and is
not contradicted by this family.

Conjecture~5.3 of that preprint states that every general rank sequence is
eventually polynomial and every particular rank sequence is eventually
``pseudo-polynomial.'' The rank-thirteen example has coefficient cancellation,
so it concerns the particular case. The term ``pseudo-polynomial'' is not
defined in the cited text. Theorem~\ref{thm:main} proves failure of eventual
quasipolynomiality in the explicit sense defined here; it does not depend on
identifying these two terms and is not asserted to refute Conjecture~5.3
without that interpretation. No refutation of the general-rank assertion or
Conjecture~6.2 is claimed. The citation identifies the author-hosted text;
its correspondence with a specific arXiv version, subsequent corrections,
and historical priority remain separate questions.

\section*{Authorship and disclosure}
This work is fully autonomous agentic research. GPT 6.1 Sol in Codex
produced formal proofs and verification software for the accompanying
package. Dev D. Goyal is responsible for publication.

\bibliographystyle{plain}
\bibliography{references}
\end{document}
